\chapter{光的折射}\label{ch:g10:refraction}

把吸管插进一杯水中：它在水面处看起来像被折断了一样。
游泳池总比看上去的要深，钻石能从一块无色的石头中闪出彩色的
火光，而你正在读的这句话，
很可能借由一根比头发还细的玻璃丝穿越了大洋。一条定律
就能解释以上四种现象——两个正弦之间的等式，由斯涅耳于
1621 年发现，在此之前，物理学家盯着角度表格看了十五个世纪，
却始终没有看出其中的规律。

\section{反射：回顾}

在均匀透明的介质中，光沿直线传播——
即\cref{ch:g10:light-spectra}中的光线。两种介质的分界面处，
一部分光被弹回（\emph{反射}），另一部分
穿过界面（\emph{\omterm{def:g10:refraction:refraction}{折射}}）。两者都遵循精确的几何定律，
并采用同一约定：角度总是从\emph{\omterm{def:g10:refraction:angles}{法线}}量起，
而不是从表面量起。

\begin{definition}[法线与入射角]\label{def:g10:refraction:angles}
光线与表面相交的点处，
\emph{法线}\index{法线}是该点处与表面
垂直的直线。\emph{入射角}\index{入射角} $i_1$ 是入射光线
与法线之间的夹角；反射角和\omterm{def:g10:refraction:refraction}{折射}角
的测量方式与之相同，
均从法线量起。
\end{definition}

\begin{proposition}[反射定律]\label{prop:g10:refraction:reflection}
反射光线位于由入射光线与
\omterm{def:g10:refraction:angles}{法线}所确定的平面内，位于\omterm{def:g10:refraction:angles}{法线}另一侧，且反射角
等于\omterm{def:g10:refraction:angles}{入射角}：
\[
i_1' = i_1 .
\]
\end{proposition}
\admitted

\begin{example}[正确读取角度]\label{ex:g10:refraction:mirror}
一束激光“以 $35^\circ$”射到平面镜上——该角度是从
镜面量起的，光束通常也是这样描述的。\omterm{def:g10:refraction:angles}{入射角}为
$90^\circ - 35^\circ = 55^\circ$，因此反射光线也与\omterm{def:g10:refraction:angles}{法线}成
$55^\circ$ 角射出，光束转过的角度为
$180^\circ - 2 \times 55^\circ = 70^\circ$。忘记
\omterm{def:g10:refraction:angles}{法线}约定是本章最常见的错误，它每次都会
使结果偏差 $90^\circ$。
\end{example}

\section{折射与折射率}

把铅笔插入半满的玻璃杯中：它在水面处看起来像是断了。
水下部分发出的光离开水面时方向会改变，
因此眼睛——它默认光线沿直线传播——会把笔尖
定位到实际并不在的位置。

\begin{definition}[折射]\label{def:g10:refraction:refraction}
\emph{折射}\index{折射}是指光穿过两种透明介质的
分界面时传播方向的改变。第二种介质中的
光线称为\emph{折射光线}\index{折射光线}，其与\omterm{def:g10:refraction:angles}{法线}的夹角
即折射角 $i_2$。
\end{definition}

\omterm{def:g10:refraction:refraction}{折射}之所以发生，是因为光在不同介质中的
速率不同；每种材料的特征在于它使光减速的程度。

\begin{definition}[折射率]\label{def:g10:refraction:index}
透明介质的\emph{折射率}\index{折射率}
是以下比值
\[
n = \frac{c}{v},
\]
其中 $c = \qty{3.00e8}{m/s}$ 为真空中的光速，$v$ 为光
在该介质中的速率。由于 $v \leq c$，\omterm{def:g10:refraction:refraction}{折射}率满足 $n \geq 1$；它
没有\omterm{def:g10:orders-of-magnitude:unit}{单位}。一些常用的数值如下：
\begin{center}
\begin{tabular}{lccccc}
介质 & 空气 & 水 & 有机玻璃 & 玻璃 & 金刚石 \\
$n$ & $1.0003$ & $1.33$ & $1.49$ & $1.50$ & $2.42$ \\
\end{tabular}
\end{center}
\end{definition}

\begin{example}[光在水中的速率]\label{ex:g10:refraction:speed}
在水中，$v = \dfrac{c}{n} = \dfrac{\qty{3.00e8}{m/s}}{1.33}
= \qty{2.26e8}{m/s}$。在金刚石中，光的速率仅为
$\qty{3.00e8}{m/s}/2.42 = \qty{1.24e8}{m/s}$——不到其真空中速率的
一半。对于空气，$n = 1.0003$ 使速率改变 $0.03\%$：
本章中将空气视为真空，取 $n_{\text{空气}} = 1.00$。
\end{example}

\begin{theorem}[斯涅耳折射定律]\label{thm:g10:refraction:snell}
\index{斯涅耳定律}
\omterm{def:g10:refraction:refraction}{折射光线}位于由入射光线和\omterm{def:g10:refraction:angles}{法线}
所确定的平面内，位于\omterm{def:g10:refraction:angles}{法线}另一侧，且\omterm{def:g10:refraction:angles}{入射角}与
\omterm{def:g10:refraction:refraction}{折射}角满足
\[
n_1 \sin i_1 = n_2 \sin i_2 ,
\]
其中 $n_1$ 和 $n_2$ 分别为两种介质的\omterm{def:g10:refraction:index}{折射率}。光
路是可逆的：光沿\omterm{def:g10:refraction:refraction}{折射光线}逆向传播时，
将沿入射光线射出。
\end{theorem}
\admitted

\begin{remark}\label{rem:g10:refraction:honest}
此处将反射定律与\omterm{def:g10:refraction:refraction}{折射}定律陈述为实验
事实。它们在本科一年级分册中得到严格推导，那里把光
视为波：两条定律——以及公式 $n = c/v$ 本身——
则都可由波在光密介质中减速而得出。
\end{remark}

\begin{omfigure}
\begin{tikzpicture}[scale=1.0, font=\small]
  \fill[blue!8] (-3.2,-2.1) rectangle (3.2,0);
  \draw[thick] (-3.2,0) -- (3.2,0);
  \draw[dashed] (0,-2.1) -- (0,2.3);
  \coordinate (O) at (0,0);
  \coordinate (I) at (-1.8,2.0);
  \coordinate (R) at (1.8,2.0);
  \coordinate (T) at (1.0,-1.73);
  \coordinate (N) at (0,2.3);
  \coordinate (Nd) at (0,-2.1);
  \draw[-Stealth, omDef, thick] (I) -- (O);
  \draw[-Stealth, omExo, thick] (O) -- (R);
  \draw[-Stealth, omThm, thick] (O) -- (T);
  \pic[draw, angle radius=9mm, "$i_1$", angle eccentricity=1.35] {angle = N--O--I};
  \pic[draw, angle radius=7mm, "$i_1'$", angle eccentricity=1.4] {angle = R--O--N};
  \pic[draw, angle radius=9mm, "$i_2$", angle eccentricity=1.35] {angle = Nd--O--T};
  \node[anchor=west] at (-3.1,0.45) {介质 1（$n_1$）};
  \node[anchor=west] at (-3.1,-0.55) {介质 2（$n_2$）};
\end{tikzpicture}

{\small 入射光线（蓝色）、反射光线（灰色）与\omterm{def:g10:refraction:refraction}{折射光线}（红色）。角度均自\omterm{def:g10:refraction:angles}{法线}（虚线）量起，
光线进入较密介质（$n_2 > n_1$）时，
\emph{朝向}\omterm{def:g10:refraction:angles}{法线}偏折。}
\end{omfigure}

\begin{example}[空气到水]\label{ex:g10:refraction:airwater}
一束光线以 $i_1 = 30^\circ$ 的\omterm{def:g10:refraction:angles}{入射角}
射到平静的池面上。取 $n_1 = 1.00$，$n_2 = 1.33$：
\[
\sin i_2 = \frac{n_1 \sin i_1}{n_2} = \frac{\sin 30^\circ}{1.33}
= \frac{0.500}{1.33} = 0.376,
\qquad
i_2 = 22.1^\circ .
\]
光线向\omterm{def:g10:refraction:angles}{法线}偏折，这总是发生在进入较慢介质
（$n_2 > n_1$）时。离开水面时，同样的计算倒过来进行，
光线则\emph{远离}\omterm{def:g10:refraction:angles}{法线}偏折。
\end{example}

\begin{method}[解折射问题的方法]\label{met:g10:refraction:method}
\begin{enumerate}
  \item 画出界面、\omterm{def:g10:refraction:angles}{法线}和光线，确定 $n_1$
        （入射光线所在的介质）与 $n_2$。
  \item 检查所给角度是否自\omterm{def:g10:refraction:angles}{法线}量起；
        若自界面量起，则需换算。
  \item 写出 $n_1 \sin i_1 = n_2 \sin i_2$，分离出未知的
        正弦值；再用计算器的反正弦功能求出角度，
        与数学分册中的做法相同。
  \item 合理性检查：若 $n_2 > n_1$，光线向\omterm{def:g10:refraction:angles}{法线}偏折，
        若 $n_2 < n_1$，则远离\omterm{def:g10:refraction:angles}{法线}偏折，且 $\sin i_2$ 的结果
        必须介于 $0$ 与 $1$ 之间。
\end{enumerate}
\end{method}

定律背后的实验只需一张桌子即可完成：一块有机玻璃半圆盘，
一把量角器、一台激光器。原始的角度数据对 $(i_1, i_2)$ 看起来杂乱无章——
托勒密约在公元 150 年将其制成表格，并错误地断定
$i_2$ 与 $i_1$ 成正比。若改为以正弦对正弦作图，则
数据会排成一条过原点、斜率为 $n_1/n_2$ 的直线：这条
直线\emph{即是}斯涅耳定律。

\begin{omfigure}
\begin{tikzpicture}
\begin{axis}[omaxis, width=7.2cm, height=5.4cm,
  xlabel={$\sin i_1$}, ylabel={$\sin i_2$},
  xmin=0, xmax=1.08, ymin=0, ymax=0.88,
  xtick={0,0.2,0.4,0.6,0.8,1.0}, ytick={0.2,0.4,0.6,0.8}]
  \addplot[omThm, thick, domain=0:1.02] {x/1.33};
  \addplot[only marks, mark=*, mark size=1.6pt, omDef] coordinates
    {(0.174,0.131) (0.342,0.257) (0.500,0.376) (0.643,0.483)
     (0.766,0.576) (0.866,0.651) (0.940,0.707)};
\end{axis}
\end{tikzpicture}

{\small 空气到水的测量，$i_1 = 10^\circ$ 至 $70^\circ$：
以正弦对正弦作图，数据落在一条过原点的
斜率为 $1/1.33$ 的直线上。}
\end{omfigure}

\begin{proposition}[视深]\label{prop:g10:refraction:depth}
\index{视深}
在水下深度为 $d$ 处的物体，从上方以小角度观察，
看起来位于\emph{视深}处：
\[
d' = \frac{d}{n} ,
\]
其中 $n$ 为水的\omterm{def:g10:refraction:index}{折射率}。深度为 \qty{1.20}{m} 的池塘，看上去
只有 $1.20/1.33 = \qty{0.90}{m}$ 深——这是造成
鲁莽跳水的典型原因。
\end{proposition}
\admitted

\begin{remark}\label{rem:g10:refraction:depthhonest}
该公式由斯涅耳定律应用于射入眼睛的近似竖直光线而得；
推导过程见本科一年级分册。效应的方向无需公式即可判断：
离开水面的光线向远离\omterm{def:g10:refraction:angles}{法线}的方向偏折，
眼睛沿其反向追溯时，会定位到一点，
该点高于物体的实际位置。
\end{remark}

\section{色散：每种颜色一条定律}

三棱镜能将白色阳光变为彩虹，正如
\cref{ch:g10:light-spectra}中的\omterm{def:g10:light-spectra:white}{光谱}所示。\omterm{def:g10:refraction:refraction}{折射}解释了\emph{原因}：
介质的\omterm{def:g10:refraction:index}{折射率}对不同颜色并不完全相同。

\begin{definition}[色散]\label{def:g10:refraction:dispersion}
称介质为\emph{色散的}\index{色散}，当其
\omterm{def:g10:refraction:index}{折射率}随光的\omterm{def:g10:light-spectra:wavelength}{波长}而变。在玻璃和水中，
$n$ 对蓝光（短\omterm{def:g10:light-spectra:wavelength}{波长}）比对红光（长\omterm{def:g10:light-spectra:wavelength}{波长}）略大，
因此蓝光偏折得更多。
\end{definition}

\begin{example}[两种颜色，一个表面]\label{ex:g10:refraction:twocolors}
一束\omterm{def:g10:light-spectra:white}{白光}以 $i_1 = 60^\circ$ 射到冕玻璃上。对于这种
玻璃，$n_{\text{红}} = 1.51$，$n_{\text{蓝}} = 1.53$。斯涅耳定律
对每种颜色分别应用：
\[
\sin i_{\text{红}} = \frac{\sin 60^\circ}{1.51} = 0.574,
\quad i_{\text{红}} = 35.0^\circ ;
\qquad
\sin i_{\text{蓝}} = \frac{\sin 60^\circ}{1.53} = 0.566,
\quad i_{\text{蓝}} = 34.5^\circ .
\]
单个表面使\omterm{def:g10:light-spectra:white}{白光}分开的角度仅约半度——在窗玻璃上几乎
看不见，但三棱镜两次\omterm{def:g10:refraction:refraction}{折射}，偏折方向相同，
扇面展宽；相距几米处，各色光之间已相隔数厘米。
\end{example}

\begin{remark}\label{rem:g10:refraction:rainbow}
彩虹的计算与此相同，只是在雨滴内部进行：阳光
射入水滴时发生\omterm{def:g10:refraction:refraction}{折射}，在水滴后壁反射，离开时再次\omterm{def:g10:refraction:refraction}{折射}。
水的\omterm{def:g10:refraction:index}{折射率}从红光的 $1.331$ 变化到蓝光的 $1.343$，
因此每种颜色以各自的角度射出——没有两位观察者
看到的彩虹完全相同。
\end{remark}

\section{全反射}

斯涅耳定律中暗藏一个陷阱。光从水射向空气时，\omterm{def:g10:refraction:refraction}{折射
光线}偏离\omterm{def:g10:refraction:angles}{法线}：$\sin i_2 = 1.33 \sin i_1 > \sin i_1$。
增大 $i_1$，公式很快就要求 $\sin i_2 > 1$，这是
任何角度都无法满足的。此时\omterm{def:g10:refraction:refraction}{折射}便不再发生。

\begin{definition}[临界角与全反射]\label{def:g10:refraction:critical}
设光在\omterm{def:g10:refraction:index}{折射率}为 $n_1$ 的介质中射向\omterm{def:g10:refraction:index}{折射率}较小的
介质 $n_2 < n_1$。称\emph{临界角}\index{临界角}
$i_c$ 为使\omterm{def:g10:refraction:refraction}{折射光线}恰好掠过界面的\omterm{def:g10:refraction:angles}{入射角}，
（$i_2 = 90^\circ$）。当 $i_1 > i_c$ 时，不存在\omterm{def:g10:refraction:refraction}{折射光线}：界面
将所有光按照反射定律全部反射回第一种介质，
这种现象称为\emph{全
反射}\index{全反射}。
\end{definition}

\begin{proposition}[临界角公式]\label{prop:g10:refraction:criticalformula}
当 $n_1 > n_2$ 时，
\[
\sin i_c = \frac{n_2}{n_1} .
\]
\end{proposition}

\begin{proof}
在斯涅耳定律中取 $i_2 = 90^\circ$：$n_1 \sin i_c = n_2 \sin 90^\circ
= n_2$。当 $i_1 > i_c$ 时，斯涅耳定律要求
$\sin i_2 = \frac{n_1}{n_2}\sin i_1 > 1$：不存在\omterm{def:g10:refraction:refraction}{折射}方向。
（光此时被\emph{全部}反射——而非被吸收——
是一个实验事实，由本科一年级分册中的
波动理论证实。）
\end{proof}

\begin{omfigure}
\begin{tikzpicture}[scale=0.92, font=\small]
  \fill[blue!8] (-3.6,-1.9) rectangle (3.6,0);
  \draw[thick] (-3.6,0) -- (3.6,0);
  \draw[dashed] (-2.2,-1.6) -- (-2.2,1.6);
  \draw[dashed] (0.3,-1.6) -- (0.3,1.6);
  \draw[dashed] (2.2,-1.6) -- (2.2,1.6);
  \draw[-Stealth, omDef, thick] (-2.95,-1.30) -- (-2.2,0);
  \draw[-Stealth, omThm, thick] (-2.2,0) -- (-1.20,1.12);
  \draw[-Stealth, omDef, thick] (-0.83,-0.99) -- (0.3,0);
  \draw[-Stealth, omThm, thick] (0.3,0) -- (1.75,0.06);
  \draw[-Stealth, omDef, thick] (0.90,-0.75) -- (2.2,0);
  \draw[-Stealth, omDef, thick] (2.2,0) -- (3.50,-0.75);
  \node at (-2.2,-1.85) {$i_1 < i_c$};
  \node at (0.3,-1.85) {$i_1 = i_c$};
  \node at (2.2,-1.85) {$i_1 > i_c$};
\end{tikzpicture}

{\small 光从水射向空气，\omterm{def:g10:refraction:angles}{入射角}逐渐增大：先是\omterm{def:g10:refraction:refraction}{折射光线}偏离\omterm{def:g10:refraction:angles}{法线}，
然后是\omterm{def:g10:refraction:refraction}{折射光线}恰好掠过界面（\omterm{def:g10:refraction:critical}{临界角}），再后是\omterm{def:g10:refraction:critical}{全反射}——
界面此时已成为理想的
镜子。}
\end{omfigure}

\begin{example}[三个临界角]\label{ex:g10:refraction:criticalvalues}
射向空气时：对于水，$\sin i_c = 1/1.33 = 0.752$，故
$i_c = 48.8^\circ$；对于玻璃，$\sin i_c = 1/1.50$，故
$i_c = 41.8^\circ$；对于金刚石，$\sin i_c = 1/2.42 = 0.413$，故
$i_c = 24.4^\circ$。\omterm{def:g10:refraction:index}{折射率}越大，逸出锥角越小：
在金刚石中，任何与刻面\omterm{def:g10:refraction:angles}{法线}夹角大于 $24.4^\circ$ 的光线
都会被困住——这正是其闪光的秘密，分析见
\cref{pb:g10:refraction:1}。
\end{example}

\section{光纤}

\omterm{def:g10:refraction:critical}{全反射}使一根玻璃丝变成了光的管道：
几乎平行于轴线射入的光线，射到管壁时的\omterm{def:g10:refraction:angles}{入射角}远大于
\omterm{def:g10:refraction:critical}{临界角}，无损耗地反射，每千米可反射
数百万次而不泄漏。

\begin{definition}[光纤]\label{def:g10:refraction:fiber}
\emph{光纤}\index{光纤}是一根细玻璃丝，由
\omterm{def:g10:refraction:index}{折射率}为 $n_1$ 的\emph{纤芯}\index{纤芯}外包
\omterm{def:g10:refraction:index}{折射率}略小的\emph{包层}\index{包层}（$n_2 < n_1$）构成。
从一端射入的光沿纤芯传播，依靠多次\omterm{def:g10:refraction:critical}{全
反射}，\omterm{def:g10:refraction:critical}{全反射}发生在纤芯与包层的交界面上。
\end{definition}

\begin{omfigure}
\begin{tikzpicture}[scale=1.0, font=\small]
  \draw[omExo, thick] (0,0.85) -- (7,0.85);
  \draw[omExo, thick] (0,-0.85) -- (7,-0.85);
  \fill[blue!8] (0,-0.5) rectangle (7,0.5);
  \draw[thick] (0,0.5) -- (7,0.5);
  \draw[thick] (0,-0.5) -- (7,-0.5);
  \draw[-Stealth, omThm, thick] (0,-0.36) -- (1.0,0.5) -- (2.16,-0.5)
    -- (3.32,0.5) -- (4.48,-0.5) -- (5.64,0.5) -- (6.6,-0.33);
  \node[anchor=south] at (3.5,0.9) {包层 ($n_2 = 1.46$)};
  \node[fill=blue!8, inner sep=2pt] at (3.5,-0.25) {纤芯 ($n_1 = 1.48$)};
\end{tikzpicture}

{\small 光借助\omterm{def:g10:refraction:critical}{全反射}沿\omterm{def:g10:refraction:fiber}{光纤}的\omterm{def:g10:refraction:fiber}{纤芯}被导引
传播（图中反射角被大大夸大了：实际光线几乎
以与\omterm{def:g10:refraction:angles}{法线}夹角超过 $80^\circ$ 的方向掠射壁面）。}
\end{omfigure}

\begin{example}[用数据看光纤]\label{ex:g10:refraction:fibernumbers}
设 $n_1 = 1.48$，$n_2 = 1.46$，
\[
\sin i_c = \frac{1.46}{1.48} = 0.986,
\qquad i_c = 80.6^\circ :
\]
只有与轴线夹角约在 $9^\circ$ 以内的光线才能被导引——这
没有问题，因为这正是激光入射的方式。在\omterm{def:g10:refraction:fiber}{纤芯}中，光
的传播速率为 $v = c/1.48 = \qty{2.03e8}{m/s}$：一个脉冲
穿过长为 \qty{100}{km} 的\omterm{def:g10:refraction:fiber}{光纤}约需半毫秒。
\end{example}

\begin{remark}[数据传输的数量级]\label{rem:g10:refraction:datarates}
海底电缆——即铺设于海底、由几十根\omterm{def:g10:refraction:fiber}{光纤}组成的
束——承担了几乎全部洲际通信流量。单根现代\omterm{def:g10:refraction:fiber}{光纤}
通过同时发送许多\omterm{def:g10:light-spectra:wavelength}{波长}，其传输速率的\omterm{def:g10:orders-of-magnitude:oom}{数量级}为 $10^{13}$ 比特每秒；
整条电缆的传输速率可达 $10^{14}$ 比特每秒，
大约每隔 \qty{80}{km} 设一个放大器。相比之下，卫星
承担的流量仅占百分之一的极小部分。
\end{remark}

\section{练习}

\begin{exercise}[$\star$]\label{exo:g10:refraction:1}
一束激光射到平面镜上，入射方向与\emph{镜面
所成的}角为 $25^\circ$。求\omterm{def:g10:refraction:angles}{入射角}、反射角，以及入射光束与
反射光束之间的夹角。
\end{exercise}

\begin{exercise}[$\star$]\label{exo:g10:refraction:2}
\begin{enumerate}
  \item 求光在水中（$n = 1.33$）和金刚石
        中（$n = 2.42$）的传播速率。
  \item 某固体中光的传播速率为 \qty{1.97e8}{m/s}，求其
        \omterm{def:g10:refraction:index}{折射率}，并判断该材料是何种物质。
\end{enumerate}
\end{exercise}

\begin{exercise}[$\star$]\label{exo:g10:refraction:3}
一束光从空气射到玻璃砖（$n = 1.50$）上，\omterm{def:g10:refraction:angles}{入射角} $i_1 = 40^\circ$。
求\omterm{def:g10:refraction:refraction}{折射}角。玻璃内的光线应以多大的\omterm{def:g10:refraction:angles}{入射角}射到表面，
才能以
$40^\circ$ 的\omterm{def:g10:refraction:refraction}{折射}角射入空气？
\end{exercise}

\begin{exercise}[$\star$]\label{exo:g10:refraction:4}
阳光以与竖直方向成 $55^\circ$ 的角射到湖面。
求\omterm{def:g10:refraction:refraction}{折射光线}的角度，并说明光线是偏向\omterm{def:g10:refraction:angles}{法线}
还是偏离\omterm{def:g10:refraction:angles}{法线}——并用一个词说明理由。
\end{exercise}

\begin{exercise}[$\star$]\label{exo:g10:refraction:5}
为鉴别一种透明液体，将一束光射到其表面，
\omterm{def:g10:refraction:angles}{入射角} $i_1 = 45^\circ$，测得\omterm{def:g10:refraction:refraction}{折射}角 $i_2 = 32^\circ$。
求该液体的\omterm{def:g10:refraction:index}{折射率}，并对照
\cref{def:g10:refraction:index}中的表格，判断其为何种液体。
\end{exercise}

\begin{exercise}[$\star\star$]\label{exo:g10:refraction:6}
一块透明塑料半圆盘给出以下测量数据：
\begin{center}
\begin{tabular}{lccccc}
$i_1$ & $20^\circ$ & $30^\circ$ & $40^\circ$ & $50^\circ$ & $60^\circ$ \\
$i_2$ & $13.5^\circ$ & $19.5^\circ$ & $25.5^\circ$ & $31.0^\circ$
      & $35.5^\circ$ \\
\end{tabular}
\end{center}
\begin{enumerate}
  \item 对每组数据计算 $\dfrac{\sin i_1}{\sin i_2}$，
        观察到什么现象？
  \item 由此推出该塑料的\omterm{def:g10:refraction:index}{折射率}，并判断其为何种材料。
  \item 为什么以 $\sin i_2$ 对 $\sin i_1$ 作图，
        比以 $i_2$ 对 $i_1$ 作图更能检验斯涅耳定律？
\end{enumerate}
\end{exercise}

\begin{exercise}[$\star\star$]\label{exo:g10:refraction:7}
求水、玻璃（$n = 1.50$）和金刚石相对于空气的
\omterm{def:g10:refraction:critical}{临界角}。再用斯涅耳定律说明，为何光线\emph{从}空气
\emph{射入}水中时没有\omterm{def:g10:refraction:critical}{临界角}。
\end{exercise}

\begin{exercise}[$\star\star$]\label{exo:g10:refraction:8}
\begin{enumerate}
  \item 水池深 \qty{2.0}{m}，站在池边的游泳者
        看到的水深为多少？
  \item 一只苍鹭注视着水面下 \qty{40}{cm} 处游动的鱼。
        鱼看起来位于多深处？苍鹭应
        瞄准像的上方、像处，还是像的下方？
\end{enumerate}
\end{exercise}

\begin{exercise}[$\star\star$]\label{exo:g10:refraction:9}
一束光射到一块窗玻璃（$n = 1.50$，两表面平行）上，\omterm{def:g10:refraction:angles}{入射角}为 $50^\circ$。
\begin{enumerate}
  \item 求光在第一个表面的\omterm{def:g10:refraction:refraction}{折射}角。
  \item \omterm{def:g10:refraction:refraction}{折射光线}从内侧以相同的
        角射到第二个表面，求出射角，并将其与 $50^\circ$ 比较。
  \item 那么，厚窗玻璃对穿过它的光线的方向与
        位置各有什么影响？
\end{enumerate}
\end{exercise}

\begin{exercise}[$\star\star$]\label{exo:g10:refraction:10}
\omterm{def:g10:light-spectra:white}{白光}以 $55^\circ$ 的\omterm{def:g10:refraction:angles}{入射角}射到一块火石玻璃上。对于这种
玻璃，$n_{\text{红}} = 1.61$，$n_{\text{蓝}} = 1.66$。求红光和
蓝光的\omterm{def:g10:refraction:refraction}{折射}角，以及它们之间的夹角。
哪种颜色的光更接近\omterm{def:g10:refraction:angles}{法线}？
\end{exercise}

\begin{exercise}[$\star\star$]\label{exo:g10:refraction:11}
一条\omterm{def:g10:refraction:fiber}{光纤}的\omterm{def:g10:refraction:fiber}{纤芯}\omterm{def:g10:refraction:index}{折射率}为 $1.48$，
\omterm{def:g10:refraction:fiber}{包层}\omterm{def:g10:refraction:index}{折射率}为 $1.46$。
\begin{enumerate}
  \item 求\omterm{def:g10:refraction:fiber}{纤芯}与\omterm{def:g10:refraction:fiber}{包层}交界面上的\omterm{def:g10:refraction:critical}{临界角}。
  \item 一条导光线恰好以\omterm{def:g10:refraction:critical}{临界角}沿\omterm{def:g10:refraction:fiber}{光纤}来回反射前进。证明
        在 \qty{1.0}{km} 长的直\omterm{def:g10:refraction:fiber}{光纤}中，它比沿轴线直行的
        光线多走约 \qty{14}{m}。
\end{enumerate}
\end{exercise}

\begin{exercise}[$\star\star\star$]\label{exo:g10:refraction:12}
一位潜水者的眼睛位于水面以下 \qty{2.0}{m} 处，水面完全平静。由于
\omterm{def:g10:refraction:refraction}{折射}，潜水者看到\emph{整个}天空被压缩成
正上方的一个明亮圆盘——斯涅耳窗。
\begin{enumerate}
  \item 解释为何来自地平线的光线以水的\omterm{def:g10:refraction:critical}{临界角}到达潜水者的眼睛，
        其与竖直方向的夹角为 $48.8^\circ$。
  \item 求水面上明亮圆盘的半径，利用
        由眼睛、竖直线与圆盘边缘构成的直角三角形中的正切。
  \item 在水面上，潜水者在圆盘\emph{之外}看到什么？
\end{enumerate}
\end{exercise}

\begin{exercise}[$\star\star\star$]\label{exo:g10:refraction:13}
一个棱镜的三个角分别为 $30^\circ$、$60^\circ$、$90^\circ$。一细束
\omterm{def:g10:light-spectra:white}{白光}垂直射入与直角相邻的一个侧面，
不偏折地穿过，并从内侧射到长侧面上，
\omterm{def:g10:refraction:angles}{入射角}为 $i_1 = 30^\circ$。这种玻璃的 $n_{\text{红}} = 1.51$ 和
$n_{\text{蓝}} = 1.53$。
\begin{enumerate}
  \item 求红光与蓝光的出射角，以及它们之间
        光束扇形的角宽度。
  \item 将棱镜换成 $45^\circ$--$45^\circ$--$90^\circ$ 棱镜，
        使内部\omterm{def:g10:refraction:angles}{入射角}变为 $45^\circ$。
        证明光线完全不会从长侧面射出。
\end{enumerate}
\end{exercise}

\begin{exercise}[$\star\star\star$]\label{exo:g10:refraction:14}
潜望镜用 $45^\circ$--$45^\circ$--$90^\circ$ 玻璃
棱镜使光线折转（$n = 1.50$）：光束垂直射入一个短侧面，
并从内侧以 $45^\circ$ 射到长侧面上。
\begin{enumerate}
  \item 证明长侧面起完美反射镜的作用。这为何
        优于金属反射镜？金属反射镜每次反射都会吸收
        百分之几的光。
  \item 潜望镜进水，长侧面现在与水接触。
        求玻璃–水交界面的新\omterm{def:g10:refraction:critical}{临界角}，
        并证明该反射镜失效。
  \item 第 (2) 问中的光线以何角度离开棱镜
        进入水中？
\end{enumerate}
\end{exercise}

\begin{exercise}[$\star\star\star$]\label{exo:g10:refraction:15}
一条跨大西洋电缆在伦敦与纽约之间铺设了
\qty{6200}{km} \omterm{def:g10:refraction:fiber}{光纤}（\omterm{def:g10:refraction:fiber}{纤芯}\omterm{def:g10:refraction:index}{折射率}为 $1.47$）。
\begin{enumerate}
  \item 求光在\omterm{def:g10:refraction:fiber}{纤芯}中的速率以及脉冲的单程
        传播时间。
  \item 无线电波以 $c$ 传播。在相同距离上，它需要多长时间，
        以及与无线电相比，\omterm{def:g10:refraction:fiber}{光纤}中的玻璃
        会多耗费多少往返时间？
  \item 一些金融公司不惜重金，只为把这条链路的时延
        缩短几毫秒。试提出两种不同的物理方法，使光
        更早到达。
\end{enumerate}
\end{exercise}

\section{问题：钻石、泳池与光纤}

\begin{problem}[{周末问题——斯涅耳定律的三种应用：一池谎报深度的泳池，
一块被切割得使光无法射出的宝石，以及
互联网之下那六十毫秒的海洋}]
\label{pb:g10:refraction:1}
同一个关于两个正弦的等式，三种用途。在游泳池中，它
让你对深度产生误判，并把整个天空压缩成一个圆盘。在
钻石中，\omterm{def:g10:refraction:angles}{入射角}被推到极限时，它便不让光射出——
几个世纪以来，珠宝匠一直围绕这一“拒绝”切割宝石。大西洋之下，
它引导激光脉冲在玻璃纤维中传播，并设定了
全球互联网的速度上限。本题依次讨论这三种应用，
每次都使用同一条定律。

\medskip
\noindent\textbf{第一部分——热身。}
\begin{enumerate}
  \item 求光在水中（$n = 1.33$）与钻石中的速率，
        钻石（$n = 2.42$）使光速减慢了多少倍？
  \item 一束在空气中的光以 $i_1 = 40^\circ$ 射到池塘上，
        求\omterm{def:g10:refraction:refraction}{折射}角。
  \item 水中一束光从下方以 $28.9^\circ$ 射到水面，
        以何角度射入空气？请说明
        这体现的是哪条原理。
  \item 求水–空气交界面的\omterm{def:g10:refraction:critical}{临界角}，
        并描述水下光线以 $60^\circ$ 射到水面时
        会发生什么。
  \item 空气的\omterm{def:g10:refraction:index}{折射率}为 $1.0003$。求光在空气中的速率，
        并说明 $n_{\text{空气}} = 1.00$ 这一近似
        在全文中使用是合理的。
\end{enumerate}

\medskip
\noindent\textbf{第二部分——泳池。}
\begin{enumerate}[resume]
  \item 泳池深 \qty{3.0}{m}。利用视深公式
        （见\cref{prop:g10:refraction:depth}），求站在池上方的人
        所感知的深度。
  \item 借助双光线示意图（以文字描述），解释为何池底的硬币
        看起来像是升高了：离开水面的光线会发生什么，
        眼睛把这些光线的交点定在何处？
  \item 一位潜水者的眼睛位于 \qty{3.0}{m} 深处。求斯涅耳窗——
        即潜水者透过它看到整个天空的明亮圆盘
        的半径（见\cref{exo:g10:refraction:12}）。
  \item 池底的灯向上发出与竖直方向成 $30^\circ$、
        $45^\circ$ 和 $50^\circ$ 的光线。对每条光线，
        说明它是否能逸出，若能，以何角度射出；若不能，则会发生什么。
  \item 岸边的叉鱼者瞄准一条鱼，叉击应瞄准鱼的像的上方、
        正对还是下方？而鱼回望时，是否也会看到
        叉鱼者同样发生了位置偏移？
\end{enumerate}

\medskip
 \noindent\textbf{第三部分——钻石。}
\begin{enumerate}[resume]
  \item 求钻石（$n = 2.42$）对空气的\omterm{def:g10:refraction:critical}{临界角}，
        并与玻璃（$n = 1.52$）的\omterm{def:g10:refraction:critical}{临界角}相比较，
        后者为 $41.1^\circ$。
  \item 宝石内部的一条光线射到某一刻面上，入射方向与其
        \omterm{def:g10:refraction:angles}{法线}成 $30^\circ$。在钻石中会发生什么？在玻璃仿制品中呢？
        用一句话解释，钻石的切工与其很小的\omterm{def:g10:refraction:critical}{临界角}相结合，
        为何会产生闪光。
  \item 钻石的\omterm{def:g10:refraction:dispersion}{色散}也很强：$n_{\text{红}} = 2.41$，
        $n_{\text{蓝}} = 2.45$。一束\omterm{def:g10:light-spectra:white}{白光}以
        $45^\circ$ 的\omterm{def:g10:refraction:angles}{入射角}射入某一刻面，求各色光的\omterm{def:g10:refraction:refraction}{折射}角及
        两者的角度差。
  \item 随后，红光与蓝光都从内部射到一个出射刻面，
        \omterm{def:g10:refraction:angles}{入射角}均为 $17.0^\circ$。求两条出射角，以及
        出射光扇的角度差，并与入射时的角度差相比较。
        \omterm{def:g10:refraction:critical}{临界角}附近，\omterm{def:g10:refraction:dispersion}{色散}会怎样变化？
  \item 一滴水（$n = 1.33$）附着在某一刻面上。求新的
        钻石–水界面的\omterm{def:g10:refraction:critical}{临界角}，判断
        第 12 题中 $30^\circ$ 光线的去向，并解释珠宝商为何
        要将钻石保持得一尘不染。
\end{enumerate}

\medskip
 \noindent\textbf{第四部分——\omterm{def:g10:refraction:fiber}{光纤}。}
\begin{enumerate}[resume]
  \item 一根\omterm{def:g10:refraction:fiber}{光纤}的\omterm{def:g10:refraction:fiber}{纤芯}\omterm{def:g10:refraction:index}{折射率}为 $1.48$，\omterm{def:g10:refraction:fiber}{包层}\omterm{def:g10:refraction:index}{折射率}为
        $1.46$。求\omterm{def:g10:refraction:fiber}{纤芯}–\omterm{def:g10:refraction:fiber}{包层}界面处的
        \omterm{def:g10:refraction:critical}{临界角}。
  \item 若将一根裸玻璃丝置于空气中，其\omterm{def:g10:refraction:critical}{临界角}为
        $42.5^\circ$——看似好得多。请给出仍要采用\omterm{def:g10:refraction:fiber}{包层}的两个理由。
        （可以想一想灰尘、划痕，以及表面被接触会对
        \omterm{def:g10:refraction:critical}{全反射}产生什么影响。）
  \item 求光在\omterm{def:g10:refraction:fiber}{纤芯}中的速率，再求一个脉冲
        沿 \qty{6000}{km} \omterm{def:g10:refraction:fiber}{光纤}传播的单程时间。
  \item 以恰好等于\omterm{def:g10:refraction:critical}{临界角}的角度反复\omterm{def:g10:refraction:critical}{全反射}的光线，
        其行进的路程是\omterm{def:g10:refraction:fiber}{光纤}轴线长度的 $1/\sin i_c$ 倍。求
        比 \qty{6000}{km} 多出的路程，以及由此多花的时间。
  \item 压轴的一问。一对现代\omterm{def:g10:refraction:fiber}{光纤}每秒可传输约 $10^{13}$ 比特，
        而本书约含 $2.4 \times 10^8$ 比特。
        求横跨大洋的往返（“ping”）时间，以及
        \omterm{def:g10:refraction:fiber}{光纤}每秒传送的本书册数——
        然后用一句话说明斯涅耳定律
        对互联网有何作用。
\end{enumerate}
\end{problem}
